\documentclass[11pt,a4paper]{ctexart}

\usepackage[margin=2.2cm]{geometry}
\usepackage{amsmath,amssymb}
\usepackage[dvipsnames]{xcolor}
\usepackage{hyperref}
\usepackage{enumitem}
\usepackage[most]{tcolorbox}

\hypersetup{colorlinks=true,linkcolor=blue!55!black,urlcolor=blue!55!black}
\setlist[enumerate]{itemsep=0.35em,topsep=0.45em}
\setlist[itemize]{itemsep=0.35em,topsep=0.45em}

\tcbset{
  commonbox/.style={enhanced,breakable,boxrule=0.8pt,arc=2mm,
    left=1.4mm,right=1.4mm,top=1mm,bottom=1mm,colback=white}
}
\newtcolorbox{coreblock}{commonbox,colframe=BrickRed!65!black,colback=BrickRed!4,
  title=核心论断,fonttitle=\bfseries}
\newtcolorbox{quoteblock}{commonbox,colframe=RoyalBlue!70!black,colback=RoyalBlue!4,
  title=原文引述,fonttitle=\bfseries}
\newtcolorbox{caseblock}{commonbox,colframe=Orange!70!black,colback=Orange!5,
  title=典型偏差,fonttitle=\bfseries}
\newtcolorbox{actionblock}{commonbox,colframe=ForestGreen!60!black,colback=ForestGreen!4,
  title=怎么做,fonttitle=\bfseries}
\newtcolorbox{keypointblock}{commonbox,colframe=black!65,colback=black!3,
  title=一句话记住,fonttitle=\bfseries}
\newtcolorbox{exampleblock}[1]{commonbox,colframe=Purple!60!black,colback=Purple!4,
  title=#1,fonttitle=\bfseries}

\title{怎样坚持对上负责和对下负责相统一\\
  \large 《求是》2026/08 张忠军 文章解读}
\author{}
\date{\today}

\begin{document}
\maketitle

\begin{coreblock}
对上负责和对下负责，从来都是\textbf{统一的、不可分割的}：对党负责就是对人民
负责，对人民负责就是对党负责，二者统一于\textbf{对党和人民事业的高度负责}
之中。这是树立和践行\textbf{正确政绩观}的内在要求。
\end{coreblock}

\tableofcontents
\newpage

\section{文章框架一览}

文章分三个层次推进：

\begin{enumerate}[label=\textbf{\arabic*.}]
  \item \textbf{是什么}：界定``对上负责''与``对下负责''的内涵，论证二者
        在党性与人民性上的内在一致。
  \item \textbf{反什么}：剖析现实中把二者对立或割裂的几种典型偏差，
        及其根源。
  \item \textbf{怎么办}：从思想、行动、制度、用人四个维度提出落实办法。
\end{enumerate}

\section{一、对上、对下负责到底是什么}

\subsection*{对上负责}

\begin{quoteblock}
``地方和部门的权威都来自于党中央权威''；地方和部门的工作``都是对党中央
决策部署的具体落实''。
\end{quoteblock}

\textbf{对上负责}=对上级党组织负责，归根到底=对党中央负责。
要求：

\begin{itemize}
  \item 始终坚持党中央集中统一领导；
  \item 把``两个维护''落到实处；
  \item 不折不扣贯彻党中央决策部署。
\end{itemize}

\begin{exampleblock}{生活例子：公司里的``对上负责''}
公司总部要求各分公司``今年要把客户投诉率降到 1\% 以下''。
对上负责的分公司经理会认真研究总部为什么提这个要求（提升品牌口碑、
减少售后成本），然后制定可落地的方案，定期向总部汇报进展。
\textbf{偏差版}：经理只盯着``投诉率''数字，把投诉电话转接给外包、
把差评删掉，数字达标了，但客户实际体验更差——这是把``对上负责''
异化成``对数字负责''，而不是对总部背后的真实目标负责。
\end{exampleblock}

\subsection*{对下负责}

\textbf{对下负责}=对人民群众负责，也包括对下一层级的党组织和党员、干部负责。
要求两条腿走路：

\begin{itemize}
  \item \textbf{面向群众}：以人民为中心，把好事实事做到群众心坎上；
  \item \textbf{面向下级}：上级党组织要切实关心、支持下级，为基层减负、为
        担当者撑腰。
\end{itemize}

\begin{exampleblock}{生活例子：社区书记的``对下负责''}
社区要推进老旧小区加装电梯。对下负责的书记会逐户走访，
了解一楼住户担心采光、顶楼住户担心分摊费用，然后协调出``一楼减免、
政府补贴、高层分摊''的方案，让大多数居民真正受益。
\textbf{偏差版}：书记只听``热心居民''（往往是几个活跃分子）的意见，
开一次会就拍板，结果大多数沉默居民的利益被忽视——
这是把``对下负责''简化成``对嗓门大的人负责''。
\end{exampleblock}

\subsection*{为什么二者本来就是一回事}

\begin{quoteblock}
毛泽东同志在党的七大政治报告中指出：``全心全意地为人民服务，一刻也不脱离
群众；一切从人民的利益出发，而不是从个人或小集团的利益出发''。
\end{quoteblock}

党的根本宗旨是\emph{为人民服务}，党的事业\emph{就是}人民的事业。所以：

\[
\text{对党负责} \;\Longleftrightarrow\; \text{对人民负责}.
\]

把任何一方与另一方对立起来，都是党性出了问题。

\begin{keypointblock}
对上对下不是两本账，而是同一本账的两个面：党性即人民性。
\end{keypointblock}

\section{二、被反对的几种典型偏差}

文章用四个``画像'' 描出干部最容易踩的坑。

\begin{caseblock}{}
\textbf{偏差 1：唯上哲学}
\begin{itemize}
  \item 把对上负责异化成对上级领导个人负责；
  \item 挖空心思迎合领导，热衷于打造领导``可视范围''内的项目；
  \item ``不怕群众不满意，就怕领导不注意''。
\end{itemize}
\end{caseblock}

\begin{exampleblock}{生活例子：``盆景工程''}
某县为了迎接上级检查，把通往县城的主干道两旁刷成统一白墙，
却对背街小巷的垃圾堆积视而不见。领导坐车路过时``眼前一亮''，
但老百姓出门还是绕道走。这就是典型的``不怕群众不满意，
就怕领导不注意''——把对上负责简化成``让领导眼前一亮''。
\end{exampleblock}

\begin{caseblock}{}
\textbf{偏差 2：片面对下负责（本位主义）}
\begin{itemize}
  \item 只顾一域不管全局；
  \item 对党中央决策做选择、搞变通、打折扣；``下有对策''软抵制``上有政策''；
  \item 重要问题先斩后奏、边斩边奏，甚至斩而不奏。
\end{itemize}
\end{caseblock}

\begin{exampleblock}{生活例子：``地方保护''}
某地为保本地就业，对外地企业设置隐形门槛，不让外地建材进入本地市场。
表面上看是``对本地企业负责''，实际上破坏了全国统一大市场，
损害了全局利益。这就是把``对下负责''异化成``对本地小团体负责''，
用``下有对策''软抵制``上有政策''。
\end{exampleblock}

\begin{caseblock}{}
\textbf{偏差 3：对上糊弄、对下敷衍}
\begin{itemize}
  \item 既没吃透上情，也没摸清下情；
  \item 照本宣科，机械执行；
  \item 坐等上级下指示，不推不动，缺乏内生动力。
\end{itemize}
\end{caseblock}

\begin{exampleblock}{生活例子：``文件旅行''}
上级要求``加强基层应急能力建设''，某单位直接把文件转发给下属，
加一句``请结合实际贯彻落实''，然后就等着汇报材料。
既没研究上级文件到底要解决什么问题，也没去现场看基层缺什么、
能做什么。文件在系统里``旅行''了一圈，基层能力还是老样子。
\end{exampleblock}

\begin{caseblock}{}
\textbf{偏差 4：能力不足、政绩观偏差}
\begin{itemize}
  \item 受个人主义、功利主义、官僚主义腐蚀；
  \item 缺乏从政治上、大局上看问题的能力，吃不准上情、看不懂下情。
\end{itemize}
\end{caseblock}

\begin{exampleblock}{生活例子：``数字政绩''}
某乡镇为了完成``新增市场主体 500 家''的指标，动员干部挨家挨户
劝居民注册个体户，甚至把一户拆成``餐饮''``零售''两家来凑数。
数字上去了，但大多数``市场主体''根本没有实际经营。
这是功利主义政绩观的典型：只盯着上级考核的数字，不管实际效果，
也不问群众是否需要这些``市场主体''。
\end{exampleblock}

\subsection*{偏差的根子}

\begin{enumerate}[label=\textbf{\arabic*.}]
  \item \textbf{主观}：党性不强，政绩观跑偏。
  \item \textbf{能力}：政治站位不高，调查研究不深。
  \item \textbf{制度}：考核评价不科学，用人导向不到位。
\end{enumerate}

\section{三、怎样把上下真正打通}

文章给出四条抓手：\textbf{思想、行动、责任、用人}。

\subsection{接好``天线''，铺好``地线''}

\begin{quoteblock}
习近平总书记强调：``任何事情都要向上看看，向下看看。''
\end{quoteblock}

\begin{actionblock}
\begin{itemize}
  \item \textbf{接天线}：吃透党中央精神、上级部署的方向性原则。
  \item \textbf{铺地线}：把这些原则结合本地实际、回应群众关切，落到具体工作。
\end{itemize}
\end{actionblock}

\begin{exampleblock}{生活例子：``双减''政策的落地}
中央要求``减轻义务教育阶段学生作业负担和校外培训负担''。
接好天线的教育局，会先吃透：这不是``不让补课''，而是``让教育回归学校主阵地''。
铺好地线后，他们不是简单关停所有培训机构，而是：
\begin{itemize}
  \item 在学校内增加课后服务、兴趣课程；
  \item 对合规机构发许可证、对违规机构依法整治；
  \item 向家长宣讲政策初衷，减少焦虑。
\end{itemize}
结果：学生负担真降了，家长焦虑也缓解了——这才是上下打通。
\end{exampleblock}

\subsection{正确对待组织}

\begin{quoteblock}
党员在党组织面前``不能隐瞒自己、信口雌黄''；说老实话、办老实事、做老实人。
\end{quoteblock}

\begin{actionblock}
\begin{itemize}
  \item 上级决定要不讲价钱、不打折扣地执行；
  \item 有不同意见，按组织程序提出，但不得违令而行；
  \item 这是\emph{对党忠诚度}的直接检验。
\end{itemize}
\end{actionblock}

\begin{exampleblock}{生活例子：``坏消息也要报''}
某县发生山洪，道路中断。基层干部面临选择：
\begin{itemize}
  \item \textbf{对上糊弄版}：怕被问责，先压着不报，等雨停了再说``已妥善处置''；
  \item \textbf{对下敷衍版}：只发个通知让``注意安全''，不组织转移；
  \item \textbf{正确版}：第一时间按程序上报，同时组织群众转移，
        哪怕要承担``前期预警不足''的责任。
\end{itemize}
正确对待组织，就是``坏消息也要报''，因为隐瞒只会让上级误判、
让群众遭更大灾。
\end{exampleblock}

\subsection{结合实际，创造性落实}

\begin{quoteblock}
毛泽东同志：``盲目地表面上完全无异议地执行上级的指示\dots\ 是反对上级
指示或者对上级指示怠工的最妙方法。''
\end{quoteblock}

党中央的部署多为\textbf{原则性、方向性}的，不可能细化到一地一域。所以：

\begin{actionblock}
\begin{itemize}
  \item 既不能机械照搬，也不能借``结合实际''之名打折扣；
  \item 要在大政方针之下，做精做实落实方案。
\end{itemize}
\end{actionblock}

\begin{exampleblock}{生活例子：垃圾分类的``本地化''}
中央要求``普遍推行垃圾分类''，这是方向性原则。
\begin{itemize}
  \item \textbf{机械照搬版}：不管本地垃圾处理能力如何，直接照搬上海标准，
        居民分了半天，最后还是混装混运——群众怨声载道。
  \item \textbf{结合实际版}：先调研本地处理设施，分阶段推进：
        第一阶段只分``可回收/其他''，同时建分拣中心；
        第二阶段再细化到``厨余/有害''等四分类。
\end{itemize}
创造性落实，就是在大方向不变的前提下，找到本地群众能接受、
设施能跟上的节奏。
\end{exampleblock}

\subsection{敢于担当，守土尽责}

\begin{actionblock}
\begin{itemize}
  \item 守土有责、守土负责、守土尽责；
  \item 不向上交应担之责，不向下推该扛之事；
  \item 推动责任落实加强督查，防止``开会等于落实、布置等于完成''。
\end{itemize}
\end{actionblock}

\begin{exampleblock}{生活例子：``踢皮球''与``接烫手山芋''}
某小区外墙脱落，居民找社区，社区说``归街道管''；
街道说``归住建局管''；住建局说``要申请专项资金''。
三家都``对上负责''（按流程走），但没人``对下负责''（居民头顶悬着砖）。
\textbf{担当版}：街道主动牵头，先拉警戒线、做临时加固（马上对下负责），
同时向上级申请资金、协调住建局（同步对上负责），
而不是把``烫手山芋''踢来踢去。
\end{exampleblock}

\subsection{用好考核评价与选人用人导向}

\begin{actionblock}
\begin{itemize}
  \item 鲜明树立\textbf{重实干、重实绩、重担当}的用人导向；
  \item 大力选拔政治过硬、敢于担当、锐意改革、清正廉洁的干部；
  \item 健全政绩考核评价体系，干部考核\textbf{充分听取群众意见}；
  \item 对做不到上下统一的，坚决处理、果断调整。
\end{itemize}
\end{actionblock}

\begin{exampleblock}{生活例子：``群众打分''进考核}
某地在干部年度考核中，拿出 40\% 的权重让服务对象（群众、企业）
打分，而不是只看上级评价。
结果：那些``领导面前表现好、群众面前不作为''的干部排名靠后，
而那些``默默修路通水、群众口碑好''的干部被提拔。
这就是用考核指挥棒把``对上''和``对下''真正绑在一起。
\end{exampleblock}

\section{把全篇拧成一根绳}

\begin{enumerate}[label=\textbf{\arabic*.}]
  \item \textbf{是什么}：对上=对党，对下=对人民；二者统一于党的事业。
  \item \textbf{反什么}：唯上、本位、糊弄、能力不足。
  \item \textbf{怎么办}：接天线 + 铺地线，老实对组织，结合实际抓落实，
        担当尽责，并用考核与用人导向兜底。
\end{enumerate}

\begin{keypointblock}
判断一个干部上下是否真打通，看的不是表态、不是开会、不是文件，
而是\textbf{经得起实践、人民、历史检验的实绩}。
\end{keypointblock}

\section{自检与延伸思考}

读完文章，可用下面几个问题自检：

\begin{enumerate}[label=\textbf{Q\arabic*.}]
  \item 我手上正在推进的工作，向上是否对得起党中央的部署方向？
        向下是否真正回应了群众关切？
  \item 我做工作时，更怕``领导不满意''，还是更怕``群众不满意''？
        是否两边都怕、都在意？
  \item 落实上级要求时，是机械搬运、是搞变通，还是结合本地实际做了
        创造性方案？
  \item 遇到棘手问题，我是先斩后奏、还是按组织程序请示报告？
  \item 在我管的下属面前，我是只让他们听话，还是给他们撑腰、为他们减负？
\end{enumerate}

\bigskip
\noindent\textbf{原文出处：}张忠军，《怎样坚持对上负责和对下负责相统一》，
《求是》2026 年第 8 期，
\url{https://www.qstheory.cn/20260415/cc79e5e1a8fb458497ed057023306336/c.html}。

\end{document}
